\documentclass[leqno]{jsarticle}
\usepackage{amssymb}
\usepackage{amsthm}
\usepackage{amsmath}
\usepackage{comment}
%定理環境 thmはあまり使わずpropをなるべく使う。
%主張は基本的に証明の中で使い、そのため番号を付けない。
%注意環境を付けてみた。
\newtheorem{thm}{定理}
\newtheorem{prop}[thm]{命題}
\newtheorem{cor}[thm]{系}
\newtheorem{lem}[thm]{補題}
\newtheorem{df}[thm]{定義}
\newtheorem{exam}[thm]{例}
\newtheorem{fact}[thm]{事実}
\newtheorem*{claim}{主張}
\newtheorem*{cau}{注意}
\renewcommand{\proofname}{{\bf 証明}}
%矢印たち。
%\toは\rightarrowと同じ。\getsは\leftarrowと同じ。同値は\iff
%上向き矢はuparrow逆はdownarrow
\newcommand{\To}{\Longrightarrow}
\newcommand{\Gets}{\LongLeftarrow}
%黒板太字
\newcommand{\nn}{\mathbb{N}}
\newcommand{\zz}{\mathbb{Z}}
\newcommand{\qq}{\mathbb{Q}}
\newcommand{\rr}{\mathbb{R}}
\newcommand{\cc}{\mathbb{C}}
%無限大
\newcommand{\mgn}{\infty}
%差集合は\setminusで統一する。等号の否定は\neq
%真部分集合は\subsetneqq　偏微分は\partial
%ディスプレイスタイル。\dfracでディスプレイ分数
\newcommand{\dpst}{\displaystyle}
\newcommand{\ep}{\varepsilon}

%空集合は\Oを使う！！！！！！！！！！！！

\title{ストーンチェックコンパクト化の構成}
\author{一色三良(concious77)}
\date{\today}
			\begin{document}
\maketitle
この文書では$T_{3.5}$空間のストーンチェックコンパクト化を構成する。構成するだけなので細かい性質の証明は省略する。
\begin{df}[完全正則空間]
点$p$とその点を含まない閉集合$C$との任意の組$(p,C)$が関数で分離出来る。つまり任意の組$(p,C)$に対してある実連続関数$f:X\to [0,1]$で
\[
p\in f^{-1}(0),\ C\subset f^{-1}(1)
\]
を充たすものが存在するとき、その空間を{\bf 完全正則}であると言う。また、完全正則でかつ$T_1$であるような空間を$T_{3\frac{1}{2}}$空間もしくは$T_{3.5}$空間と呼ぶ。
\end{df}
上と同じ記号を用いると、	
$p\in f^{-1}([0,\frac{1}{2}),f^{-1}([0,\frac{1}{2})\cap C=\O$で$f^{-1}([0,\frac{1}{2})$は$X$の開集合なので完全正則空間は正則空間である。よって$T_{3.5}$空間は$T_3$空間である。

\begin{prop}
$\mathcal{C}$を完全正則空間$(X,\mathfrak{O})$から単位閉区間$[0,1]$への連続関数全体とする。このとき、完全正則空間の位相は$\mathcal{C}$から定まる始位相
$(X,\mathfrak{I}(\mathcal{C}))$\footnote{この文書では関数族$\mathcal{A}$から誘導される始位相を$\mathfrak{I}(\mathfrak{A})$と書くことにする。}と一致する。\footnote{実はこの命題の逆も成立するがすこし面倒なので証明はしない。\cite{2}参照のこと}

\end{prop}
\begin{proof}
$\mathcal{C}$の元は$(X,\mathfrak{O})$の連続関数なので明らかに
\[
\mathfrak{I}(\mathcal{C})\subset \mathfrak{O}
\]
が成立する。逆向きの包含関係を示そう。$U\in \mathfrak{O}$で$x\in U$とすると、$X\setminus U$は$(X,\mathfrak{O})$の閉集合で$p$を含まないので$X$の完全正則性から連続関数$f:X\to [0,1]$で
\[
p\in f^{-1}(0),\ X\setminus U\subset f^{-1}(1)
\]
となるものが存在する。$\mathcal{C}$の定義から明らかに$f\in \mathcal{C}$で
\[
p\in f^{-1}([0,\frac{1}{2}))\subset U
\]
なので$f^{-1}([0,\frac{1}{2}))\in \mathfrak{I}(\mathcal{C})$から
\[
\mathfrak{O}\subset \mathfrak{I}(\mathcal{C})
\]
\end{proof}


次の始位相に関する埋め込み定理を思い出そう。
\begin{fact}[始位相の埋め込み定理]\label{main}位相空間$(X,\mathfrak{O})$には$\mathcal{F}=\{f_{\lambda}\}_{\lambda \in \Lambda}$から定まる始位相が入っているとしこの関数族の終域を$\{(Y_{\lambda},\mathfrak{O}_{\lambda})\}_{\lambda \in \Lambda}$とするこのとき関数
\[
\phi:X\rightarrow \prod_{\lambda\in \Lambda}Y_{\lambda \in \Lambda}\ ;x\mapsto(f_{\lambda}(x))_{\lambda\in \Lambda}
\]について$\mathfrak{I}(\mathcal{F})=\mathfrak{I}(\phi)$が成立する。
また$\phi$が単射ならばこれは埋め込みになっているつまり$(X,\mathfrak{O})$は$\prod_{\lambda\in \Lambda}Y_{\lambda \in \Lambda}$の部分空間$\phi(X)$と同相となる。
\end{fact}
証明は電波通信の記事にあると思う。

この埋め込み定理と先の命題から次のことが直ちに従う。記号は上のものと同じとする。

\begin{thm}
$T_{3.5}$空間$X$は
\[
\phi:X\to\prod_{f\in \mathcal{C}}[0,1]  ;x\mapsto (f(x))_{f\in \mathcal{C}}
\]
によって
$\dpst \prod_{f\in \mathcal{C}}[0,1]=[0,1]^{\mathcal{C}}=I^{\mathcal{C}}$に埋め込める。
\end{thm}
\begin{df}
$T_{3.5}$空間$X$のストーンチェックコンパクト化$\beta X$を
\[
\beta X=CL_{I^{\mathcal{C}}}(\phi (X))
\]
と定義する。
\end{df}


\section{継ぎ足し}
\begin{df}[ゼロ集合、コゼロ集合]
位相空間$X$とその上の連続関数$f:X\to \rr$について$f$の{\bf ゼロ集合}$Z(f)$を
\[
Z(f)=f^{-1}(0)
\]
と定義し、{\bf コゼロ集合}$C(f)$を
\[
C(f)=X\setminus Z(f)=\{x\in X\ |\ f(x)\neq 0\}
\]
と定義する。明らかに$Z(f)$は$X$の閉集合で$C(f)$は$X$の開集合である。また明らかにゼロ集合の補集合はコゼロ集合で、コゼロ集合の補集合はゼロ集合である。
\end{df}

\begin{prop}$X,Y$を位相空間、
$f:X\to Y$を連続写像、$A\subset Y$を$Y$のゼロ集合とする。このとき、$f^{-1}(U)$は$X$のゼロ集合である。
\end{prop}
\begin{proof}
$g:Y\to \rr$を$Z(g)=A$となる連続関数とする。すると
\[
Z(g\circ f)=(g\circ f)^{-1}(0)=f^{-1}(g^{-1}(0))=f^{-1}(A)
\]
なので定理は明らかである。
\end{proof}
双対的に次の命題が成り立つ。
\begin{prop}$X,Y$を位相空間、
$f:X\to Y$を連続写像、$U\subset Y$を$Y$のコゼロ集合とする。このとき、$f^{-1}(U)$は$X$のコゼロ集合である。
\end{prop}
この命題と$\rr$の開集合はすべてコゼロ集合であることから次の系が得られる。
\begin{cor}
位相空間$X$と連続関数$f:X\to \rr$に及び$\rr$の開集合$U$について$f^{-1}(U)$はコゼロ集合である。
\end{cor}


\begin{prop}
位相位相$X$のゼロ集合$A$について、ある連続関数$g:X\to [0,1]$が存在して
\[
A=Z(g)
\]
が成り立つ。
\end{prop}


\begin{prop}
位相位相$X$のコゼロ集合$U$について、ある連続関数$g:X\to [0,1]$が存在して
\[
U=C(g)
\]
が成り立つ。
\end{prop}



\begin{prop}
位相空間$X$上の２つの連続関数$f,g:X\to [0,+\mgn)$について
\[
Z(f)\cap Z(g)=Z(\max(f,g)),\ Z(f)\cup Z(g)=Z(\min(f,g))
\]
が成り立つ。双対的に
\[
C(f)\cup C(g)=C(\max(f,g)),\ C(f)\cap C(g)=C(\min(f,g))
\]
が成り立つ。
\end{prop}
\begin{proof}
明らか
\end{proof}

\begin{cor}
ゼロ集合全体の族は有限個の共通部分及び、有限個の和集合で閉じている。\par
双対的にコゼロ集合全体の族も有限個の共通部分及び、有限個の和集合で閉じている
\end{cor}


\begin{cor}
$\mathfrak{I}(\mathcal{C})=\mathfrak{O}$となる位相空間$(X,\mathfrak{O})$は完全正則空間
\end{cor}


\begin{thebibliography}{9}
\bibitem{1}森田紀一,位相空間論,岩波全書331,岩波書店1981
\bibitem{2}M.G.Murdeshwar,general topology,Wiley Eastern Limited,1983
\end{thebibliography}




			\end{document}



\begin{comment}
[\bigin \endを使わない方法]

{\prop[aa] aaa}
で
\begin{prop}[aa]
aaa
\end{prop}
と同じ\df \thm \proof でも同様

[直和]
$\amalg$

[同値関係でよく使う波線]
\sim

[引数付きの新命令]
\newcommand{\prn}[1]{\|#1\|_p}

[番号のリセット]

\setcounter{equation}{0}


[字体の変更]

{\rm hogehoge}と使う。

\rm ローマン
\it イタリック
\sl 斜体

[supp]

\newcommand{\supp}{\text{{\rm supp}}}

[中央揃えの改行]

	\begin{eqnarray*}
		hoge1&=&hogehoge
		hoge2&=&hogehofe

	\end{eqnarray*}

&&の部分で揃えられる。






[場合分け]

	\begin{cases}
		hoge1 & \text{状況１}\\
		hoge2 & \text{状況２}\\
		・
		・
		・
	\end{cases}



\end{comment}





